\section{Results}
\label{sec:results}

\subsection{RQ1: Quality on ISD-Agent-Bench}
\label{sec:rq1}

\Cref{tab:rq1} reports the seven agents on \texttt{test\_90}. AlignmentGraph-ISD attains the highest composite and ADDIE rubric score under
\emph{both} models (87.77 / 85.46 on the LLM; 86.47 / 83.96 on the SLM). Wilcoxon
signed-rank tests paired by scenario, with Holm correction over the six
comparisons, place it above every baseline in both runs with
$p_{\mathrm{Holm}} < 1.4\times10^{-13}$ in all twelve comparisons; the mean paired
margin ranges from $+2.20$ (vs.\ Baseline) to $+8.31$ (vs.\ EduPlanner) on the LLM,
and from $+3.59$ (vs.\ Dick-Carey-Agent) to $+30.46$ (vs.\ EduPlanner) on the SLM,
with bootstrap 95\% CIs excluding zero throughout. The advantage is concentrated in
the artifact rubric rather than the trajectory score, consistent with the harness
improving \emph{what} is produced rather than how the tool trace reads.

\begin{table}[t]
  \caption{RQ1: ISD-Agent-Bench \texttt{test\_90} (two-judge means;
  Total = composite, ADDIE = artifact rubric, Traj.\ = trajectory;
  \textbf{bold} = best per column).}
  \label{tab:rq1}
  \small
  \begin{tabular}{@{}lrrr@{\hskip 10pt}rrr@{}}
    \toprule
    & \multicolumn{3}{c}{DeepSeek-v4-flash (LLM)} & \multicolumn{3}{c}{Qwen3.5-9B (SLM)} \\
    \cmidrule(lr){2-4}\cmidrule(lr){5-7}
    Agent & Total & ADDIE & Traj. & Total & ADDIE & Traj. \\
    \midrule
    Baseline          & 85.56 & 84.24 & 88.64 & 80.81 & 77.97 & 87.44 \\
    EduPlanner        & 79.47 & 76.16 & 87.20 & 56.01 & 43.49 & 85.23 \\
    ADDIE-Agent       & 83.12 & 79.45 & 91.68 & 80.89 & 76.35 & 91.49 \\
    RPISD-Agent       & 80.98 & 79.81 & 83.71 & 79.32 & 77.05 & 84.61 \\
    Dick-Carey-Agent  & 84.71 & 82.28 & 90.38 & 82.88 & 79.78 & 90.14 \\
    ReAct-ADDIE       & 83.37 & 79.20 & \textbf{93.08} & 80.08 & 74.50 & \textbf{93.10} \\
    \textbf{AlignmentGraph-ISD} & \textbf{87.77} & \textbf{85.46} & \textbf{93.16} & \textbf{86.47} & \textbf{83.96} & 92.34 \\
    \bottomrule
  \end{tabular}
\end{table}

\paragraph{Model scale.} The harness is nearly invariant to scale: its composite
falls only $1.30$ points from the LLM to the SLM ($87.77 \rightarrow 86.47$), while
the baselines degrade far more---Baseline loses $4.75$ ($85.56 \rightarrow 80.81$)
and EduPlanner collapses by $23.46$ ($79.47 \rightarrow 56.01$). Because the harness
holds its quality while the model weakens, its paired margin over the strongest
baseline \emph{widens} from $+2.20$ on the LLM (over Baseline) to $+3.59$ on the SLM
(over Dick-Carey-Agent). This is the harness-over-model effect made concrete: the
advantage comes from the execution environment, so it grows exactly where model
capability shrinks.

\subsection{RQ2: Ablation}
\label{sec:rq2}

\Cref{tab:rq2} isolates each mechanism on both models. Removing multi-agent
decomposition (\texttt{wo-ma}) is by far the most damaging: $+3.90$ composite and
$+3.29$ ADDIE points on the LLM ($p_{\mathrm{Holm}} = 2.3\times10^{-13}$), and
\emph{twice} as large on the SLM at $+7.94$ composite and $+7.79$ ADDIE
($p_{\mathrm{Holm}} = 7.0\times10^{-14}$; 95\% CI $[+5.46, +10.85]$). The
contribution of graph-based decomposition thus \emph{grows as the model shrinks},
mirroring the harness-over-model effect of RQ1. Removing per-agent
verification/repair (\texttt{wo-qc}) changes the composite by only $+0.26$ (LLM) and
$+0.15$ (SLM) points, in both cases not statistically distinguishable from zero
($p_{\mathrm{Holm}} = 0.89$ and $0.53$).

\begin{table}[t]
  \caption{RQ2 ablation (composite Total; contribution = \texttt{full} $-$ arm,
  paired). The decomposition contribution roughly doubles on the SLM.}
  \label{tab:rq2}
  \small
  \begin{tabular}{@{}lrr@{}}
    \toprule
    & DeepSeek (LLM) & Qwen3.5-9B (SLM) \\
    \midrule
    \texttt{full} (complete harness)   & 87.77 & 86.47 \\
    \texttt{wo-ma} ($-$ decomposition) & 83.86 & 78.53 \\
    \texttt{wo-qc} ($-$ verify/repair) & 87.51 & 86.33 \\
    \midrule
    \multicolumn{3}{@{}l}{\emph{Contribution (paired composite; $p_{\mathrm{Holm}}$):}}\\
    \quad decomposition & $+3.90$ ($2.3{\times}10^{-13}$) & $+7.94$ ($7.0{\times}10^{-14}$) \\
    \quad verify/repair & $+0.26$ (0.89, n.s.) & $+0.15$ (0.53, n.s.) \\
    \bottomrule
  \end{tabular}
\end{table}

\paragraph{Does verify/repair do anything?} That its judge-score contribution is
zero does not mean it is inactive. It fires on the high-value checks---Bloom-verb
validity, rubric consistency, active-practice coverage, and assessment--objective
cognitive-level match---and the residual audit shows it measurably reduces alignment
defects. On the SLM, the \texttt{wo-qc} arm leaves $39$ residual content-level
violations across the 90 outputs, whereas \texttt{full} leaves only $12$---a
$69\%$ reduction, concentrated in exactly the checks the tools repair (Bloom
validity, rubric consistency, cognitive-level match). Thus verification/repair does
its job; that reduction is simply not rewarded by the rubric---a construct gap we
return to in \cref{sec:discussion}. (The offline audit reconstructs the graph from
output text without edges, so its absolute per-type counts are approximate; the
across-arm comparison and the in-run repair activity are the reliable signals.)
